
\begin{ledgroupsized}[r]{120mm}
\footnotesize
\pstart
\noindent\textbf{Überlieferung:}
\pend
\end{ledgroupsized}
\begin{ledgroupsized}[r]{114mm}
\footnotesize
\pstart
\parindent -6mm
\makebox[6mm][l]{\textit{L}}%
Aufzeichnungen:
LH~XXXV~8, 30~Bl.~133. 
Ein Blatt (ca.~20~x~16~cm.);
{Wasserzeichenfragment} am Rand (französisches Papier);
alle Ränder beschnitten; 
Eineinhalb Seiten auf Bl.~133~r\textsuperscript{o} bis Bl.~133~v\textsuperscript{o};
Bl. 133~r\textsuperscript{o} horizontal in fünf Felder aufgeteilt,
Bl. 133~v\textsuperscript{o} in drei Bereichen, die Aufzeichnungen auf Bl. 133~r\textsuperscript{o} in gleicher Höhe weiterführen, zur Hälfte beschrieben.
Textfolge für die erste, zweite und vierte Aufzeichnung: Bl. 133~r\textsuperscript{o}, Bl. 133~v\textsuperscript{o}; beim vierten Text ist der Seitenübergang durch den Kustos "arbitrarium" und den Hinweis "(verte)" gesichert (\refpassage{35_08_30_133_3a}{35_08_30_133_3b}).
\pend
\pstart
\noindent
Cc 2, Nr. 1421
\pend
\end{ledgroupsized}
\vspace{5mm}
\begin{ledgroup}
\footnotesize
\pstart
\noindent
\textbf{Datierungsgründe:}
Mit N.\;\ref{5konzepte133} liegt
eine Gruppe fünf auf dem gemeinsamen Träger LH~XXXV, 8, 30 Bl.~133
sequentiell verfasster Aufzeichnungen 
(\refpassage{35_08_30_133-teil1a}{35_08_30_133-teil1b}, \refpassage{35_08_30_133-teil2a}{35_08_30_133-teil2b}, \refpassage{35_08_30_133-teil3a}{35_08_30_133-teil3b}, \refpassage{35_08_30_133-teil4a}{35_08_30_133-teil4b} und \refpassage{35_08_30_133-teil5a}{35_08_30_133-teil5b})
vor,
die Leibniz ursprünglich wohl zu Streifen zerschneiden wollte.
Deshalb hat er bei drei der fünf Texte, deren Länge eine Fortsetzung auf der Rückseite erforderte, den Schlussteil auf Bl.~133~v\textsuperscript{o} in gleicher Höhe wie den Anfang geschrieben. (Eine ähnliche Textverteilung durch sequentielle Abfassung liegt bei dem Träger LH~XXXV, 8, 30 Bl.~121, der Notizen aus der Pariser Zeit überliefert; zwei von diesen, N.\;\ref{Cc2_847-1} und N.\;\ref{Cc2_847-2}, werden als Nachträge im vorliegenden Band ediert.)
Wegen der konsekutiven Beschreibung des Blattes ist anzunehmen, dass Leibniz die Aufzeichnungen, oder zumindest ihre Anfänge auf Bl.~133~r\textsuperscript{o}, sequentiell verfasst hat.
Gemäß des Charakters des Blattes als einer Sammlung thematisch verwandter Aufzeichnungen beleuchten alle fünf Texte den Begriff Schwerpunkt: Der erste und dritte behandeln die Bestimmung des Schwerpunkts von Rotationsfiguren; der zweite und vierte die Bestimmung des gemeinsamen Schwerpunkt von Punkten in eigentümlichen Konfigurationen bzw. Bewegungen; der fünfte die einem Körper eigene Bewegung in Zusammenhang mit der des gemeinsamen Schwerpunkts des Systems.
\pend
\pstart
Zur Datierung von N.\;\ref{Cc2-1421} lässt sich zunächst feststellen, dass das Blatt ein Fragment eines 
im Leibniz-Korpus für Februar bis Juli 1676 häufig belegten französischen Wasserzeichens enthält: 
u.\,a. in den
\cite{02169}Auszügen aus 
\indexN{\textls{Mariotte}, Edme, Seigneur de Chazeuil ca. 1620\textendash1684}Mariottes  \textit{Logique}-Handschrift (\textit{LSB} VI, 3 N.\;22, Mai bis Juli 1676)
und den Exzerpten \cite{02170}\textit{Anatomica quaedam ex manuscripto \indexN{\textls{Descartes} (Cartesius, des Cartes), René 1596\textendash1650}Cartesii} mit der editorischen Datierung Februar bis September 1676 (\textit{LSB} VII, 2 N.\;58).
Dasselbe Zeichen kommt auch in den Aufzeichnungen N.\;\ref{Cc2_1381} aus der späten Pariser Zeit (Februar bis 25. April 1676) vor, die als Nachtrag in diesem Band erscheinen.
Dennoch liefert das Wasserzeichen im Träger (im Licht noch zu erörternder inhaltlicher Umstände) keinen belastbaren Datierungsgrund: Leibniz hat die Aufzeichnungen nach dem Ende seiner Pariser Jahre verfasst.
Die Verwendung französischen Papiers erklärt sich in diesem Fall dadurch, dass Leibniz für N.\;\ref{Cc2-1421} wohl den freigebliebenen Teil eines alten Bogens bzw. Blattes verwendet hat und diesen dann abgeschnitten. Dass der Text von N.\;\ref{Cc2-1421} quer zur üblichen Beschreibrichtung des Bogens verläuft, bekräftigt die Annahme einer Zweitnutzung des französischen Papiers. (Weitere ursprünglich mit dem Träger zusammenhängende Papiere konnten bislang nicht identifiziert werden.)
\pend
\pstart
Belastbare Datierungsgründe lassen sich indessen aus inhaltlichen Betrachtungen gewinnen. Insbesondere die Aufzeichnung auf \refpassage{35_08_30_133-teil2a}{35_08_30_133-teil2b}, die Leibniz offenbar als zweite verfasst hat, lässt sich mit dem Konzept N.\;\ref{60309} produktiv vergleichen, das auf Grund des Wasserzeichens auf den Zeitraum September 1685 bis Januar 1687 datierbar ist.
Das Konzept stellt eine ausführliche Weiterentwicklung einiger Gedanken und Beweise aus der zweiten Aufzeichnung: Viele Diagramme werden wieder aufgegriffen, einige Passagen fast wörtlich wiederholt (bspw. \refnote{FnC:35_08_30_133-Lamy} hier und \refpassage{37_05_064-065_4a}{37_05_064-065_4b} in N.\;\ref{60309}).
Insbesondere die Passage \refpassage{35_08_30_133-teil2a}{35_08_30_133_4} aus der zweiten Aufzeichnung von N.\;\ref{Cc2-1421} wird in einer Randanmerkung zu \refnote{FnA:37_05_064-065_elementa-lineae} des Konzepts wortgleich übernommen.
Wegen dieses festgestellten Abhängigkeitsverhältnisses sind die Aufzeichnungen N.\;\ref{Cc2-1421} sicher vor dem Konzept entstanden; daher bietet der Endpunkt der vermutlichen Entstehungsspanne von N.\;\ref{60309}, Januar 1687, einen Terminus ante quem für N.\;\ref{Cc2-1421}.
\pend
\pstart
Eine Anspielung in derselben Aufzeichnung bietet wiederum einen sicheren Terminus post quem für ihre Entstehung, damit für die der gesamten Gruppe: Denn Leibniz hat den zweiten Text sehr wohl kurz nach dem ersten verfasst, wie der Schreibduktus nahelegt.
Auf \refnote{FnC:35_08_30_133-Lamy} erwähnt Leibniz eine These
\indexN{\textls{Lamy} (L'Amy, Lami), Bernard (Bernhard) 1640\textendash1715}Bernard Lamys zu Schwerpunkt und Rektifikation einer Kurve.
Die Anspielung konnte nicht mit Sicherheit aufgelöst werden: Sie gilt eventuell einer Erwähnung der Problematik in
\textsc{B.\;Lamy}, \cite{01476}\textit{Traitez de mecanique}, Paris 1679: erstes Traktat ("De l'equilibre des solides"), S.\;17--152, bes. S.\;71\,f., wobei die Stelle allerdings sehr vage und beiläufig ist.
Dennoch steht fest, dass die \textit{Traitez} von 1679 das früheste mehrerer Werke Lamys auf dem naturwissenschaftlichen Gebiet sind: Daher bietet die Anspielung einen Grund, die Entstehung der Aufzeichnungen nicht vor 1679 zu verorten.
\pend
\pstart
Eine letzte Anspielung von Interesse befindet sich in einer Randanmerkung zu \refnote{FnC:35_08_30_133-ceva} in der zweiten Aufzeichnung.
Leibniz bemerkt die Ähnlichkeit seines Ansatzes zu dem \indexN{\textls{Ceva}, Giovanni 1647\,(?)\textendash1734}Giovanni Cevas; es handelt sich um die "Appendix Geometrica" zum "De Aquis", in  \cite{02154}\textit{Opuscula mathematica}, Mailand 1682, prop.\;II, S.\;52\,f. sowie Tab.\;8, fig.\;5. (Ceva bestimmt dort allerdings den Schwerpunkts einer Halbkugeloberfläche, nicht wie Leibniz eines Halbkreises.)
Es steht fest, dass Leibniz um 1685 die "Opuscula mathematica" \cite{02153}rezensierte (\cite{01023}\textit{AE} November 1685 S. 493--496; die Rezension erscheint vsl. in einem späteren Band der Reihe), wobei er in der Besprechung den einschlägigen Anhang "Appendix geometrica" nicht erwähnte.
Dennoch wird diese Anspielung als nicht Datierungsrelevant angesehen, da es sich um einen womöglich später nachgetragenen Randzusatz handelt.
\pend
\end{ledgroup}
\selectlanguage{latin}		%
\vspace{8mm}
\pstart
\noindent
\bracks{\textit{Nachfolgender Text bis \refpassage{35_08_30_133-teil1b}{35_08_30_133-teil1b} erste Aufzeichnung:}}
\pend
\nobreak
\vspace{.5em}
\pstart
\noindent
\edlabel{35_08_30_133-teil1a}%
\href{\Scan{LH_35_08_30_133r}}%
{\lbrack133~r\textsuperscript{o}\rbrack} 
Considerandum pro ampliatione \textls{centri gravitatis},\indexS{centrum gravitatis} quid fiat si axis\indexS{axis non rectus} non sit rectus,
sed ubique eodem modo obliquus,\indexS{axis ubique eodem modo obliquus}\indexS{axis ad radium librationis obliquus} ad radium librationis\bracks{,}\indexS{radium librationis}
deinde quid si ubique variet obliquitatem.\indexS{axis variabiliter obliquus}
\pend
\begin{edfigure}%
\includegraphics[scale=1]{%
gesamttex/edit_VIII-4/images/LH_35_08_30_133_d1_133r.pdf}
\edcaption{\lbrack\textit{Fig.~1}\rbrack}
\end{edfigure}
\pstart
Ponamus axem\indexS{axis} esse cylindroidem Ellipticum\indexS{cylindroides Ellipticus} cujus sectiones sint 
\edtext{circuli circa axem\bracks{,} ita}{%
\lemma{circuli}%
\Bfootnote{%
\textit{(1)}~.~Sed
\textit{(2)}~ita ipse etsi 
\textit{(3)}~circa axem
\textit{(a)}~quod moveri possit, ipse tamen axis non potest aliter
circumagi quam circa rectam
\textit{(b)}~ita~\textit{L}}}
et res circa axem moveri possunt, et ipse axis\indexS{axis} crassus circa rectam transeuntem per centra circulorum.
Frustraque videtur hic addi axis cylindri.\indexS{axis cylindri} %
\pend
\pstart
Quid de axe curvilineo\indexS{axis curvilineus variabiliter inclinatus} ubi quovis momento
variatur inclinatio\indexS{inclinatio} ut 
\pend
\begin{edfigure}%
\includegraphics[scale=1]{%
gesamttex/edit_VIII-4/images/LH_35_08_30_133_d2_133r.pdf}
\edcaption{\lbrack\textit{Fig.~2}\rbrack}
\end{edfigure}
\pstart
Verum cum inter circumagendum statim
\edtext{asseres a se}{%
\lemma{asseres}%
\Bfootnote{%
\textit{(1)}~quasi in
\textit{(2)}~a se~\textit{L}}}
invicem divellantur inutiles sunt axes curvilinei.\indexS{axis curvilineus}
\pend
\pstart
Possumus potius concipere motum communem\indexS{motus communis} 
rectilineae perpendicularis insistentiae,
\edtext{et figuram}{%
\lemma{et}%
\Bfootnote{%
\textit{(1)}~corpus mo
\textit{(2)}~figuram~\textit{L}}}
motam superficialem\indexS{figura superficialis} vel solidam\indexS{figura solida} inter movendum
\edtext{aliquid secum}{%
\lemma{aliquid}%
\Bfootnote{%
\textit{(1)}~secum acquirere
\textit{(2)}~secum~\textit{L}}}
abripere quod prius quieverat, vel derelinquere locu et inter evolvendum intelligi potest.
\pend
\pstart
\href{\Scan{LH_35_08_30_133v}}%
{\lbrack133~v\textsuperscript{o}\rbrack}
Mirum est in omni corpore esse centrum aequilibrii\indexS{centrum aequilibrii in omni corpore existens} non centr. figurae.\indexS{centrum figurae non in omni corpore existens}
Sic omnes rectae bisecantes triangulum isosceles non transeunt per idem punctum.
Operae pretium erit inquirere in \textls{locum horum punctorum}, qui erit vicarius centri figurae.\indexS{centrum figurae}%
\edlabel{35_08_30_133-teil1b}
\pend
\vspace{.5em}
\pstart
\noindent
\lbrack\textit{Nachfolgender Text bis \refpassage{35_08_30_133-teil2b}{35_08_30_133-teil2b} zweite Aufzeichnung:}\rbrack
\pend
\nobreak
\vspace{.5em}
\pstart
\noindent
\edlabel{35_08_30_133-teil2a}%
\href{\Scan{LH_35_08_30_133r}}%
{\lbrack133~r\textsuperscript{o}\rbrack} 
Si Elementa\indexS{elementum loci centrorum gravitatis} lineae quae est locus centrorum gravitatis\indexS{locus centrorum gravitatis}\indexS{centrum gravitatis} 
sint proportionalia ipsis quorum sunt illa centra,
centrum gravitatis loci,\indexS{centrum gravitatis loci centrorum gravitatis} 
erit centrum gravitatis totale\indexS{centrum gravitatis totale} eorum quorum centra gravitatis in illum locum cadunt.%
\edlabel{35_08_30_133_4}
\pend
\begin{edfigure}%
\includegraphics[scale=1]{%
gesamttex/edit_VIII-4/images/LH_35_08_30_133_d3_133r.pdf}
\edcaption{%
\edtext{\bracks{\textit{Fig.~3}}%
}{%
\lemma{\bracks{\textit{Fig.~3}}}%
\Cfootnote[nonum]{%
Der Punkt \textit{E} hieß bei Leibniz $\beta$; diese Bezeichnung ändert Hrsg. gemäß den Angaben im Text.}}}%
\end{edfigure}
\pstart
\edtext{%
Ita locus omnium \textit{e} centrorum gravitatis\indexS{locus omnium centrorum gravitatis} triangulorum (seu sectorum infinite parvorum)
\edtext{in quae}{%
\lemma{in}%
\Bfootnote{%
\textit{(1)}~qui
\textit{(2)}~quae~\textit{L}}}
circulus vel datus sector circuli\indexS{centrum gravitatis sectoris circuli} resolvitur, est arcus circuli centro \textit{C} 
\edtext{radi\bracks{o}}{\lemma{}\Bfootnote{radio \textit{L ändert Hrsg.}}}
\textit{CE} descripti, et puncta sunt etiam centra gravitatis\indexS{centrum gravitatis} ipsorum $L\lambda$ elementorum loci centrorum.\indexS{elementum loci centrorum gravitatis}}{%
\applabel{FnC:35_08_30_133-ceva}%
\lemma{}%
\Afootnote{%
\textit{Neben dem Text und \bracks{\textit{Fig.~3}}}:
\indexN{\textls{Ceva}, Giovanni 1647\,(?)\textendash1734}Joh.~Ceva Mediol. 
sic procedit in circulo\textsuperscript{\bracks{a}}
\\[1.5ex] %
{\footnotesize 
\textsuperscript{\bracks{a}}Joh.~Ceva \bracks{...} circulo: 
Siehe \textsc{G.\;Ceva}, \cite{02154}\textit{Opuscula mathematica}, Mailand 1682: "De Aquis", "Appendix Geometrica", prop.\;II, S.\;52\,f. sowie Tab.\;8, fig.\;5. Ceva bestimmt a.\,a.\,O. allerdings den Schwerpunkts einer Halbkugeloberfläche, nicht eines Halbkreises.}}}
Et cum ratio \textit{CL} ad \textit{EB} sit 2 
\edtext{ad 3 ex}{%
\lemma{ad 3}%
\Bfootnote{%
\textit{(1)}~ob trian
\textit{(2)}~ex~\textit{L}}}
natura trianguli, hinc generaliter dici potest centrum gravitatis semicirculi\indexS{centrum gravitatis semicirculi}
vel sectoris\indexS{centrum gravitatis sectoris circuli} cujuscunque
idem esse cum centro gravitatis semicircumferentiae\indexS{centrum gravitatis semicircumferentiae} 
alterius sectoris similis et similiter positi seu concentrici
per radium subsesquialtera ratione minorem descripti.
\pend
\begin{edfigure}%
\includegraphics[scale=1]{%
gesamttex/edit_VIII-4/images/LH_35_08_30_133_d4_133r.pdf}
\edcaption{\lbrack\textit{Fig.~4}\rbrack}
\end{edfigure}
\pstart
At in Evolutione\indexS{evolutio} res ista non procedit, neque linea \textit{e}(\textit{e}) est parallela
lineae $\beta(\beta)$.
Parallela fit linea si eadem quantitas a filo semper resecetur, non si proportionalis.
Interim patet lineam centrorum gravitatis\indexS{centrum gravitatis}\indexS{linea centrorum gravitatis elementorum}
\edtext{elementorum}{\lemma{}\Bfootnote{elementorum \textit{erg.~L}}}
figurae $AC\beta A$ haberi, si semper sumatur in \textit{CB} ipsa \textit{Ce} ad \textit{CB}
ut 2 
\edlabel{35_08_30_133_1a}%
\edtext{}{%
\xxref{35_08_30_133_1a}{35_08_30_133_1b}%
\lemma{ad 3.}%
\Bfootnote{%
\textit{(1)}~Caeterum in Hemisph
\bracks{/}
\textit{(2)}~Summa
\textit{(a)}~omnium
\textit{(b)}~\textit{AC}~\textit{L}}}%
ad 3.
\pend
\begin{edfigure}%
\includegraphics[scale=1]{%
gesamttex/edit_VIII-4/images/LH_35_08_30_133_d5_133r.pdf}
\edcaption{\lbrack\textit{Fig.~5}\rbrack}
\end{edfigure}
\pstart
Summa \textit{AC}\edlabel{35_08_30_133_1b}
in $\ind{1}N\ind{2}N = NC$ in \textit{\ind{2}V\ind{1}C}.
Ergo et summae 
\edtext{aequantur ex}{\lemma{}\Bfootnote{aequantur \bars{ergo \textit{streicht Hrsg.}} ex \textit{erg.~L}}}
datis 
\edtext{summis radiorum,}{%
\lemma{summis}%
\Bfootnote{%
\textit{(1)}~\textit{AC} ad
\textit{(2)}~radiorum,~\textit{L}}}
ad axem, \textit{d} ad momentum ipsorum \textit{\ind{2}V\ind{1}C} 
\edlabel{35_08_30_133_2a}%
\edtext{}{%
\xxref{35_08_30_133_2a}{35_08_30_133_2b}%
\lemma{ex axe.}%
\Bfootnote{%
\href{\Scan{LH_35_08_30_133v}}%
{\lbrack133~v\textsuperscript{o}\rbrack} 
\textit{(1)}~Caeterum in Hemisphaerio (ut in circulo) centrum gravitatis itidem
\textit{(2)}~Similiter~\textit{L}}}%
ex axe.\indexS{axis}
\pend
\pstart
\href{\Scan{LH_35_08_30_133v}}%
{\lbrack133~v\textsuperscript{o}\rbrack} 
Similiter\edlabel{35_08_30_133_2b}
in sphaera\indexS{sphaera} vel sectore solido\indexS{sector sphaerae}\indexS{sector solidus} (ut in circulo\indexS{circulus} vel sectore plano\indexS{sector circulis}\indexS{sector planus})
centrum gravitatis\indexS{centrum gravitatis} coincidit cum centro gravitatis superficiei hemisphaerii\indexS{centrum gravitatis superficiei hemisphaerii} vel sectoris
similis et similiter positi seu concentrici, in subtertia ratione minoris, resolvendo sectorem in pyramides.\indexS{resolutio sectoris sphaerae in pyramides}
\pend
\begin{edfigure}%
\includegraphics[scale=1]{%
gesamttex/edit_VIII-4/images/LH_35_08_30_133_d6_133v.pdf}
\edcaption{\lbrack\textit{Fig.~6}\rbrack}
\end{edfigure}
\pstart
Pro curva Evoluta si ab \textit{e} ad (\textit{e}) ducantur semper 
\edtext{rectae quae}{%
\lemma{rectae}%
\Bfootnote{%
\textit{(1)}~.~Si ja
\textit{(2)}~quae~\textit{L}}}
sunt elementa loci centrorum gravitatis,\indexS{elementum loci centrorum gravitatis}\indexS{locus centrorum gravitatis} quod si jam curva
\edtext{\textit{MM}}{\lemma{}\Bfootnote{\textit{MM} \textit{erg.~L}}}
ita secetur ut ipsae \textit{e}(\textit{e}) fiant triangulis \textit{ML}(\textit{L}) proportionales\bracks{,}
idem esset centrum gravitatis curvae,\indexS{centrum gravitatis} quod 
\edtext{figurae. Verum}{%
\lemma{figurae.}%
\Bfootnote{%
\textit{(1)}~Verum
\textit{(2)}~Et ad hanc
\textit{(3)}~Verum~\textit{L}}}
revera 
\edtext{si bene ratiocinat \indexN{\textls{Lamy} (L'Amy, Lami), Bernard (Bernhard) 1640\textendash1715}Lamy,}{%
\applabel{FnC:35_08_30_133-Lamy}%
\lemma{si bene ratiocinat Lamy}%
\Cfootnote{Die Anspielung konnte nicht mit Sicherheit aufgelöst werden. Leibniz bezieht sich womöglich auf 
\textsc{B.\;Lamy}, \cite{01476}\textit{Traitez de mecanique}, Paris 1679: erstes Traktat, S.\;17--152, bes. S.\;71\,f.}}
sufficit hanc sectionem esse possibilem.
\pend
\begin{edfigure}%
\includegraphics[scale=1]{%
gesamttex/edit_VIII-4/images/LH_35_08_30_133_d7_133v.pdf}
\edcaption{\lbrack\textit{Fig.~7}\rbrack}
\end{edfigure}
\pstart
Et videtur generaliter res applicari posse ad quamlibet lineam centrorum gravitatis,\indexS{linea centrorum gravitatis quaelibet}
etiam trapeziorum\bracks{,} possumus enim semper fingere sectionem ita
\edtext{institui ut trapezia seu}{\lemma{}\Bfootnote{institui \bars{posse \textit{gestr.}} ut trapezia \bars{\textit{BC} \textit{gestr.}} seu~\textit{L}}}
rectangula \textit{B}(\textit{C}) sint ipsis \textit{e}(\textit{e}) proportionalia. 
Et coincidet centrum rectanguli cum centro\indexS{centrum cujusque elementi curvae} cujusque 
\edtext{elementi curvae.}{%
\lemma{elementi}%
\Bfootnote{%
\textit{(1)}~sed
\textit{(2)}~curvae.~\textit{L}}}
Sed jam video difficultatem\bracks{:} non incidet centrum \textit{e} rectanguli in 
medium tangentium seu rectarum per \textit{e} ductarum. Examinandum an hoc 
mutet ratiocinationem.%
\edlabel{35_08_30_133-teil2b}%
\pend
\vspace{.5em}
\pstart
\noindent
\lbrack\textit{Nachfolgender Text bis \refpassage{35_08_30_133-teil3b}{35_08_30_133-teil3b} dritte Aufzeichnung:}\rbrack
\pend
\vspace{.5em}
\nobreak
\pstart
\noindent
\edlabel{35_08_30_133-teil3a}%
\href{\Scan{LH_35_08_30_133r}}%
{\lbrack133~r\textsuperscript{o}\rbrack}
\pend
\nobreak
\begin{edfigure}[0.5em][1em]
\includegraphics[scale=1]{%
gesamttex/edit_VIII-4/images/LH_35_08_30_133_d8_133r.pdf}
\edcaption{\lbrack\textit{Fig.~8}\rbrack}
\end{edfigure}
\pstart
\noindent
\edtext{Cum movetur}{%
\lemma{Cum}%
\Bfootnote{%
\textit{(1)}~circa
\textit{(2)}~movetur~\textit{L}}}
recta per duos circulos parallelos sic ut superficiem \textls{cylindri obliqui}\indexS{cylinder obliquus} describat,
si jam recta ista intelligatur chorda alicujus arcus circuli et planum segmenti semper tendat in axem cylindri\bracks{,}\indexS{axis cylindri}
quaeritur quae superficies describatur hoc motu, quaeve lineae a quovis arcus puncto ut \textit{C}.
\pend
\pstart
Transferenda consideratio centri gravitatis\indexS{centrum gravitatis} 
a parallelis planis verticalibus corpus in aequalia momenta\indexS{plana secantia corpus in aequalia momenta} secantibus
ad convergentia in centro.\indexS{convergentia in centro gravitatis}
Si corpus sit quodcunque et in eo punctum pro arbitrio assumatur, fingatur fixum,
corpus ita se convertet sponte sua ipso manente fixo, ut quantum potest descendat\indexS{descensus quam maximum}
seu cuidam horizontali suppositae appropinquet, seu ut distantiae punctorum corporis
(potius particularum corporis, seu centrorum globulorum, vel aliorum in quae corpus resolvitur\indexS{resolutio corporis in particulas seu globulos})
ab horizonte (vel si convergunt omnia a centro gravium\indexS{centrum gravium})
sint omnium possibilium minimae.\indexS{distantiae a centro gravium omnium possibilium minimae}
Hinc si corpus \astrosun\ consideretur ut cylindricum horizontaliter jacens
\lbrack\textit{Text bricht ab.}\rbrack
\edlabel{35_08_30_133-teil3b}
\pend
\vspace{.5em}
\pstart
\noindent
\lbrack\textit{Nachfolgender Text bis \refpassage{35_08_30_133-teil4b}{35_08_30_133-teil4b} vierte Aufzeichnung:}\rbrack
\pend
\nobreak
\vspace{.5em}
\pstart
\noindent
\edlabel{35_08_30_133-teil4a}%
\href{\Scan{LH_35_08_30_133r}}%
{\lbrack133~r\textsuperscript{o}\rbrack} 
Via centri gravitatis\indexS{centrum gravitatis}\indexS{via centri gravitatis} servit ad loca Geometriae communis.\indexS{Geometria communis}
Ut si corpus unum 
\edtext{\textit{A}}{\lemma{}\Bfootnote{\textit{A} \textit{erg.~L}}}
moveatur in linea recta alterum vero \textit{B} quiescat centrum gravitatis commune\indexS{centrum gravitatis commune} 
continue movetur\indexS{motus continuus centri gravitatis communis} 
in recta priori parallela, et quia punctum \textit{C}
rectam \textit{AB} jungentem \textit{A} et \textit{B} 
in eadem semper ratione secet quae scil. est corporum,\indexS{ratio corporum}
hinc locus planum\bracks{,} si sit recta positione data \textit{A}(\textit{A}) 
et extra eam punctum \textit{B}, et \textit{AB} semper in data ratione 
\edtext{secetur\bracks{,} erit}{%
\lemma{secetur}%
\Bfootnote{%
\textit{(1)}~\textit{CC} erit
\textit{(2)}~\bars{\textit{C} \textit{streicht Hrsg.}}
\textit{(3)}~erit~\textit{L}}}
\textit{C}(\textit{C}) recta ipsi \textit{A}(\textit{A}) parallela.
\pend
\pstart
Generalius si corpora quotcunque moveantur in rectis,\indexS{corpora quotcunque mota in rectis}
aut sin aliqua ex ipsis quiescant,
via centri communis\indexS{centrum gravitatis commune} omnium\indexS{via centri communis omnium corporum motorum} erit recta\lbrack,\rbrack\
fingo enim omnia esse gravia et descendere libere\indexS{corpora gravia libere descendentia} adeoque in recta, 
utique et centrum gravitatis\indexS{centrum gravitatis} in recta descendet, nec refert etsi rectae non sint parallelae,
saltem enim descensus\indexS{descensus in plano inclinato} erit in plano inclinato.
Item etsi motus sint inaequabiles perinde est enim ac si unum in aere\indexS{corpus in aere descendens} descendat
alterum 
\edtext{in aqua;\indexS{corpus in aqua descendens}}{%
\lemma{in}%
\Bfootnote{%
\textit{(1)}~terra
\textit{(2)}~aqua;~\textit{L}}}
et unum modo in aere modo in aqua vel alio medio.\indexS{corpus in diversis mediis descendens}
Idem est si in plano horizontali vel diversis moveantur.
Si sint igitur
duae rectae positione datae, et in qualibet punctum datum, et aliud
\edlabel{35_08_30_133_3a}
\href{\Scan{LH_35_08_30_133v}}%
{\lbrack133~v\textsuperscript{o}\rbrack} 
\edtext{arbitr\bracks{ari}um,\edlabel{35_08_30_133_3b}}{\lemma{}\Bfootnote{arbitrariarium \textit{L ändert Hrsg.}}}
sic ut intervallum puncti dati et arbitrarii in una recta 
ad intervallum puncti dati et arbitrarii in alia recta sit in ratione data,
denique recta jungens duo puncta arbitraria respondentia secetur etiam in ratione data\bracks{,}
locus omnium punctorum secantium jungentes 
pro quibuslibet punctis 
\edtext{arbitrarii\bracks{s}}{%
\lemma{}%
\Bfootnote{%
arbitrarii, 
\textit{L ändert Hrsg.}}}
respondentibus erit recta.
Quid si plura puncta moveantur in superficie sphaerica\indexS{superficies sphaerica}\indexS{puncta mota in superficie sphaerica} per vias brevissimas
seu circulos magnos,\indexS{circuli magni superficies sphaericae}
an et ita eorum centrum gr.\indexS{centrum gravitatis}\bracks{?}
\pend
\pstart
Mobilia intelligi possunt puncta diversae gravitatis.\indexS{puncta diversae gravitatis}
Si plura puncta moveantur in lineis circularibus\bracks{,}\indexS{puncta mota in lineis circularibus} 
quaerenda via centri.\indexS{via centri gravitatis}
Et ita in aliis lineis dabuntur theoremata.
\pend
\pstart
Quaeratur linea centri\indexS{linea centri gravitatis} duorum mobilium quorum unum in circulo alterum in 
\edtext{recta. An}{%
\lemma{recta.}%
\Bfootnote{%
\textit{(1)}~Non videtur
\textit{(2)}~An~\textit{L}}}
generaliter summa ex ductibus\indexS{ductus mobilis in viam sui centri gravitatis} duorum mobilium in sui centri gravitatis\indexS{via centri gravitatis} vias
aequatur facto\indexS{factum ex summa mobilium in viam centri gravitatis communis} ex summa mobilium in viam centri communis.\indexS{via centri gravitatis communis}
An vero id solum verum de descensu,\indexS{descensus} vel si alias directio eadem manet.%
\edlabel{35_08_30_133-teil4b}%
\pend
\vspace{.5em}
\pstart
\noindent
\lbrack\textit{Nachfolgender Text bis \refpassage{35_08_30_133-teil5b}{35_08_30_133-teil5b} fünfte Aufzeichnung:}\rbrack
\pend
\nobreak
\vspace{.5em}
\pstart
\noindent
\edtext{\edlabel{35_08_30_133-teil5a}%
\href{\Scan{LH_35_08_30_133r}}%
{\lbrack133~r\textsuperscript{o}\rbrack} 
Sint}{%
\lemma{\href{\Scan{LH_35_08_30_133r}}%
{\lbrack133~r\textsuperscript{o}\rbrack}}%
\Bfootnote{%
\textit{(1)}~Sint
\textit{(2)}~Sit
\textit{(3)}~Sint~\textit{L}}}
duae rectae \textit{AF} et \textit{BH} positione et magnitudine datae
quae secentur in eadem ratione in punctis \textit{G} 
\edtext{et \textit{E}. Jungatur}{%
\lemma{et \textit{E}.}%
\Bfootnote{%
\textit{(1)}~Junge \textit{GE}
\textit{(2)}~Jungatur~\textit{L}}}
\textit{GE} eique parallelae ducantur \textit{AL} et \textit{FH} sic ut anguli \textit{ABL} 
\edtext{\bracks{\textit{FHC}}}{%
\lemma{}%
\Bfootnote{%
\textit{FCH}
\textit{L ändert Hrsg.}}}
sint recti.
\pend
\begin{edfigure}%
\includegraphics[scale=1]{%
gesamttex/edit_VIII-4/images/LH_35_08_30_133_d9_133r.pdf}
\edcaption{\lbrack\textit{Fig.~9}\rbrack}
\end{edfigure}
\pstart
Debent $BE \cdot FH + CE \cdot AL$ esse $= BC \cdot GE$.
Et si \textit{BC}, \textit{CE} aequales erit dupla
$GE = AL + AF$.
Usus in motibus,\indexS{motus}
possumus enim corpora\indexS{motus corporis communis et navi} considerare ut mota in navi\indexS{navis mota motu centri gravitatis} quae moveatur motu centri gravitatis\indexS{motus centri gravitatis}
et ipsis attribuere motum proprium\indexS{motus corporis proprium} ad eam rem necessarium.%
\edlabel{35_08_30_133-teil5b}%
\pend
